\chapter*{Abstract}
\addcontentsline{toc}{chapter}{Abstract}

Stable stratification suppresses vertical mixing in the atmospheric and oceanic
boundary layer, and the resulting transport is intermittent, anisotropic, and
poorly resolved by the eddy diffusivity closures that operational models still
rely on. Direct simulation of the regime is expensive because the buoyancy
Ozmidov scale collapses faster than the viscous scale as the Richardson number
grows, so a uniform grid is wasted almost everywhere. This dissertation
develops \solver, an adaptive lattice kinetic solver for the coupled momentum
and scalar problem, and uses it to characterise the turbulent Prandtl number
across the transition from continuous to intermittent turbulence.

The numerical contribution is threefold. First, a two population lattice
Boltzmann formulation is derived in which the momentum population carries a
multiple relaxation time collision operator and the scalar population carries a
single relaxation time operator on a reduced velocity set, giving a scheme whose
scalar diffusivity is independent of the momentum relaxation rate over the full
range of Prandtl numbers considered. Second, a block structured refinement
criterion is introduced that keys on the local buoyancy Ozmidov scale rather than on
a velocity gradient magnitude, so refinement follows the layers where mixing is
actually active instead of following shear that stratification has already
sterilised. Third, the coupling between refinement levels is shown to conserve
the second order moments of both populations to machine precision under a
rescaling of the non equilibrium parts, which removes the spurious entropy
production that has limited earlier refined lattice schemes at high Reynolds
number.

The solver reproduces reference statistics for unstratified channel flow at
friction Reynolds number 590 to within 1.4 percent in mean velocity and 3.1
percent in scalar variance, and reproduces the collapse of the momentum flux at
bulk Richardson number 0.25 reported by earlier spectral studies. Production
runs cover twelve stratification levels on grids of up to 1.1 billion active
lattice sites. The turbulent Prandtl number is found to rise from 0.86 in the
neutral limit to 1.9 near the critical Richardson number, and the rise is
governed by the ratio of the buoyancy Ozmidov scale to the local grid spacing rather
than by the gradient Richardson number alone, which explains the resolution
sensitivity that has confounded previous estimates. A revised closure is
proposed that carries this ratio explicitly and reduces the error in the
predicted scalar flux from 34 percent to 9 percent across the case matrix.
